% ============================================================================
% Report template - Sapienza Università di Roma
% Compile with: pdflatex report_template.tex (run twice for the TOC)
% Put the logo next to this file as  sapienza_logo.png  (or .pdf).
% If the logo is missing, a placeholder box is printed instead.
% ============================================================================
\documentclass[11pt,a4paper]{article}

\usepackage[utf8]{inputenc}
\usepackage[T1]{fontenc}
\usepackage[english]{babel}      % use [italian] for an Italian report
\usepackage{lmodern}
\usepackage[margin=2.5cm]{geometry}
\usepackage{graphicx}
\usepackage{xcolor}
\usepackage{amsmath,amssymb}
\usepackage{booktabs}
\usepackage{listings}
\usepackage{fancyhdr}
\usepackage{setspace}
\usepackage{caption}
\usepackage{tikz}
\usetikzlibrary{arrows.meta,positioning}
\usepackage[hidelinks]{hyperref}

% --- Colours ---------------------------------------------------------------
\definecolor{sapienza}{RGB}{124,31,52}     % adjust to the official maroon if needed
\definecolor{codebg}{RGB}{247,247,247}

% --- Metadata: edit only this block ----------------------------------------
\newcommand{\universityname}{Sapienza Università di Roma}
\newcommand{\facultyname}{Master's in Cybersecurity}
\newcommand{\coursename}{Statistics}
\newcommand{\reporttitle}{Shrimpgnature:\\ Signed Audit Logs under a FIFO Queue}
\newcommand{\reportsubtitle}{Design, implementation and performance analysis}
\newcommand{\studentnameG}{Guglielmo Forte}
\newcommand{\studentidG}{Matricola 2351773}
\newcommand{\studentnameL}{Staiano Luca}
\newcommand{\studentidL}{Matricola 2325802}
\newcommand{\professorname}{Prof. Tommaso Gastaldi}
\newcommand{\academicyear}{2026/2027}

% --- Helpers ---------------------------------------------------------------
\newcommand{\todo}[1]{\textcolor{red}{[TODO: #1]}}   % delete all before submitting

\lstdefinelanguage{JavaScript}{
  keywords={const,let,var,function,async,await,return,if,else,for,while,new,
            export,import,try,catch,throw,true,false,null,typeof},
  comment=[l]{//}, morecomment=[s]{/*}{*/}, morestring=[b]', morestring=[b]",
  morestring=[b]`
}
\lstset{
  language=JavaScript, basicstyle=\ttfamily\footnotesize,
  backgroundcolor=\color{codebg}, keywordstyle=\color{sapienza}\bfseries,
  commentstyle=\color{gray}\itshape, stringstyle=\color{blue!60!black},
  frame=single, rulecolor=\color{black!20}, breaklines=true,
  numbers=left, numberstyle=\tiny\color{gray}, xleftmargin=1.5em,
  framexleftmargin=1.5em, columns=fullflexible
}

% --- Header / footer -------------------------------------------------------
\pagestyle{fancy}
\fancyhf{}
\lhead{\small\textcolor{sapienza}{\universityname}}
\rhead{\small\coursename}
\cfoot{\thepage}
\renewcommand{\headrulewidth}{0.4pt}
\setlength{\headheight}{14pt}

\begin{document}

% ============================== TITLE PAGE ==================================
\begin{titlepage}
\centering

% --- LOGO LEFT ALIGNMENT START ---
\begin{flushleft}
\IfFileExists{logo_Sapienza.png}{%
    \includegraphics[width=0.33\textwidth]{./logo_Sapienza.png}%
}{\IfFileExists{sapienza_logo.pdf}{%
    \includegraphics[width=0.33\textwidth]{}%
}{%
    \fbox{\parbox[c][3.5cm][c]{5cm}{\centering\textcolor{red}{LOGO MISSING}\\ \small add \texttt{sapienza\_logo.png}}}%
}}
\end{flushleft}
\nopagebreak[4]
\vspace{1.2cm}
% --- LOGO LEFT ALIGNMENT END ---

{\Large\scshape \universityname\par}
\vspace{0.2cm}
{\large \facultyname\par}
{\large \coursename\par}
\vspace{1 cm}
{\color{sapienza}\rule{\textwidth}{1pt}}\\[0.6cm]
{\Huge\bfseries \reporttitle\par}
\vspace{0.4cm}
{\Large \reportsubtitle\par}
\vspace{0.3cm}
{\color{sapienza}\rule{\textwidth}{1pt}}
\vfill
\begin{tabular}{rl}
\textbf{Students:} & \studentnameG \\
& \studentidG \\[0.3cm]
& \studentnameL \\
& \studentidL \\[1cm]
\textbf{Professor:} & \professorname \\[3cm]
\end{tabular}
\vfill
{\large Academic Year \academicyear\par}
\end{titlepage}


% ============================== FRONT MATTER ================================
\pagenumbering{roman}
\begin{abstract}
\noindent
This report presents the ?
\end{abstract}
\tableofcontents
\clearpage
\pagenumbering{arabic}

\begin{spacing}{1.1}

% ============================== 1. INTRODUCTION (~1 page) ===================
\section{Introduction}

Digital signatures are essential for proving the authenticity, integrity and non-repudiation of documents. They guarantee that a specific record was processed by a notary and that it has not sustained any modification since being signed. \\\\
However if a system solely relied on digital signatures, a malicious actor could silently delete a previously notarized document, reorder historycal transactions or present an incomplete version of the events to an auditor which would not detect any related issues since the signatures of the individual document would still be mathematically valid.\\\\
To prevent this, a notary requires a notary requires an ordered log which enables tamper detection. By structuring recepits into a caryptographic hash chain, each new notary operation is linked to the previous one, thus creating an immutable sequence. Any attempt to remove, reorder, alter past events will inevitably cause incoherence in the chain, which is detectable by the system. Thus \textbf{digital signatures} are needed to ensure fine-grained validity of a document, while the \textbf{hash chain} operates at a coarse-grained level.

\subsection{Goals}
The main goal of this project is to simulate and analize a central notary authority that processes document signing requests. The system includes a performance simulation component that models the signing authority as a single-server queue, allowing for empirical observation of queueing metrics and comparison against theoretical models such as M/M/1 and M/G/1.

\subsection{Structure of the report}
The report is structured as follows:
\begin{itemize}
\item{\textbf{Section 2} provides background informations on digital signatures, hash chains and queueing models.}
\item{\textbf{Section 3} outlines the system architecture, describing components, message flows and data structures.}
\item {\textbf{Section 4} details the software implementation of the queue, the worker processes, the verification entity and the traffic generator.}
\item {\textbf{Section 5} presents an experimental evaluation of service and waiting times under load.}
\item {\textbf{Section 6} epitomizes the work and concludes the report.}
\end{itemize}
% ============================== 2. BACKGROUND (~2 pages) ====================
\pagebreak

\section{Background}

\subsection{Digital signatures}
Digital signatures ensure the authenticity, non-repudiation and integrity of digital documents. The central notary authority supports two asymmetric cryptographic algorithms: the Elliptic Curve Digital Signature Algorithm (\textit{ECDSA}) using the P-256 curve, and the RSA algorithm with the Probabilistic Signature Scheme (\textit{RSA-PSS}) using a 4096-bit modulus.
\\When a document digest is submitted, the authority signs it with its private key, generating a signature that is verifiable by anyone using the corresponding public key.

\subsection{Hash chains and tamper evidence}
To ensure the immutability of the \emph{audit log}, the system structures recepits into a \textit{hash chain}. Each new recepit computes a \texttt{chainHash} using the SHA-256 hash function over the previous block's hash \texttt{prevHash}, the current block's \texttt{signature} and the \texttt{docDigest}; thus forming a sequential dependency which guarantees the propagation of any tampering attempt made on any past block invalidating every subsequent one.
\\\\
The logical structure of the afore mentioned blocks can be visualized as follows:
\begin{equation}
  h_i = \mathrm{SHA\text{-}256}\bigl(h_{i-1} \,\Vert\, \sigma_i \,\Vert\, d_i\bigr),
  \qquad h_0 = 0^{256},
  \label{eq:chain}
\end{equation}
where $d_i$ is the document digest and $\sigma_i$ its signature.
\\\\
If an attacker modifies the contents of entry $j$, for example altering the document digest, the output of the hash function for that block changes entirely due to the avalanche effect, resulting in a new hash $h'_j$ \neq $h_j$. Because the next block's hash, $h_{(j+1)}$ depends on $h_j$, computing $H(h'_j \parallel \text{data}_{j+1})$ will yield a different result, $h'_{(j+1)}$. By induction this cascades forward indefinitely.

\pagebreak

\subsection{Queueing model}
The system evaluates performance using queueing theory, modelling the authority as a single-server queue with Poisson arrivals. Two analytical models are tracked to predict queue waiting times ($W_q$):

\begin{itemize}
  \item \textbf{M/M/1}: assumes both exponential inter-arrival times and exponential service times. Its infinitesimal generator matrix $Q$ is:
  \[
  Q = \begin{pmatrix}
  -\lambda & \lambda & 0 & 0 & \cdots \\
  \mu & -(\lambda + \mu) & \lambda & 0 & \cdots \\
  0 & \mu & -(\lambda + \mu) & \lambda & \cdots \\
  0 & 0 & \mu & -(\lambda + \mu) & \cdots \\
  \vdots & \vdots & \vdots & \vdots & \ddots
  \end{pmatrix}
  \]
  The expected waiting time in the queue is $W_q = \frac{\rho}{\mu (1 - \rho)}$.

  \item \textbf{M/G/1}: assumes exponential arrivals but general service times. It uses the Pollaczek-Khinchine formula: $W_q = \frac{\lambda(\sigma_S^2 + \bar{S}^2)}{2 (1-\rho)}$.
  Here, $\lambda$ is the arrival rate, $\mu$ is the service rate, $\rho = \frac{\lambda}{\mu}$ is the system load, $\bar{S}$ is the mean service time, and $\sigma_S^2$ is the service time variance.
\end{itemize}

\pagebreak

\section{System Architecture}

\subsection{Components}

The system consists of multiple modules operating asynchronously across the main browser thread and a Web Worker:

\begin{itemize}
\item \textbf{UI and Telemetry View}: handles DOM manipulations, dynamic charting with \texttt{Chart.js}, and real-time KPI updates.
\item \textbf{Traffic Generator}: produces synthetic document signing requests according to a Possion process.
\item \textbf{Central Notary Authority}: a Web Worker responisble for cryptographic key generation, document signing and hash chain computation.
\item \textbf{Verification Engine}: audits the public ledger to detect tampering or history-rewriting attacks. It also provides simulations of these attacks.
\end{itemize}


\subsection{Message flow}

When a request is generated, its payload (a document digest and algorithm choice) is enqueued in the FIFO Server. Once the server is idle, it dispatches the request via a \texttt{SIGN\_SUCCESS} message containing a receipt and telemetry data. The main thread then handles this message by updating the public audit log, telemetry metrics and the UI.

\subsection{Receipts and telemetry}
The generated receipt includes the sequence number, document digest, generated signature, chosen algorithm, timestamp and a \texttt{chainProof} (containing \texttt{prevHash} and \texttt{chainHash}). Parallel to this, telemetry data is captured for performance monitoring, tracking specifically \texttt{tArrival} (time of queue entry), \texttt{tQueueExit} (time of dispatch to the worker) and \texttt{serviceTimeSeconds} (duration of cryptographic operations).

\pagebreak{}

\section{Implementation}

\subsection{FIFO queue and single-server state machine}
The queue logic is implemented in \texttt{fifoServer.js}. It encapsulates an array-based queue and maintains the server state as either \texttt{IDLE} or \texttt{PROCESSING}. When a job is enqueued, it receives an arrival timestamp. If the server in \texttt{IDLE}, the job is immediately dequeued, its queue exit time is stamped, the state transitions to \texttt{PROCESSING}, and the request is dispatched. Upon task completion, the state resets to \texttt{IDLE}, and the \texttt{next()} routine is invoked.

\begin{lstlisting}[language=JavaScript, caption=queue/fifoServer.js]
export function createFifoServer({ dispatch, onDepthChange = () => {} }) {
  const queue = [];
  let state = "IDLE";

  function next() {
    if (state === "PROCESSING" || queue.length === 0) return;
    const request = queue.shift();
    state = "PROCESSING";
    request.tQueueExit = performance.now();
    onDepthChange(queue.length);
    dispatch(request);
  }

  function enqueue(request) {
    request.tArrival = performance.now();
    queue.push(request);
    onDepthChange(queue.length);
    next();
  }

  function complete() {
    state = "IDLE";
    next();
  }

  return { enqueue, complete, depth: () => queue.length };
}
\end{lstlisting}


\subsection{Signing and hash chaining in the worker}
The worker initializes by generating ECDSA and RSA-PSS key pairs using the native \texttt{crypto.subtle} API. During processing, the signing method is invoked over the document digest. Immediately following the signature generation, the hash chain link is computed by calling \texttt{computeHashChain}. Performance metrics strictly measure the time around these cryptographic bounds right inside the signing method. Executing this inside a Web Worker prevents CPU-intensive 4096-bit RSA operations from blocking the main thread.


\begin{lstlisting}[language=JavaScript, caption=worker/authority.js]
export function createAuthority() {
  const privateKeys = { ecdsa: null, rsa: null };
  let currentChainHash = GENESIS_HASH;
  let sequenceCounter = 0;

  async function init() {...}

  // Service time = everything from tStart to the end (sign + chain hash)
  async function sign({ id, docDigestBytes, algorithm, tArrival, tQueueExit }) {
    const tStart = performance.now();

    const signer = SIGNERS[algorithm];
    ...

    const signatureBuffer = await crypto.subtle.sign(
      signer.params,
      key,
      docDigestBytes,
    );
    const signature = bufferToHex(signatureBuffer);
    const docDigest = bufferToHex(docDigestBytes);

    const prevHash = currentChainHash;
    const chainHash = await computeChainHash(prevHash, signature, docDigest);
    currentChainHash = chainHash;
    const sequence = sequenceCounter++;

    const serviceTimeSeconds = (performance.now() - tStart) / 1000;

    return {
      receipt: {
        sequence,
        id,
        docDigest,
        signature,
        algorithm,
        timestamp: Date.now(),
        chainProof: { prevHash, chainHash },
      },
      telemetry: { tArrival, tQueueExit, serviceTimeSeconds },
    };
  }

  return { init, sign };
}
\end{lstlisting}


\subsection{Verification engine}
The audit engine (\texttt{auditLog.js}) performs a two-layer verification sequentially over the ledger. The first one involves recalculating the \texttt{chainHash} for each block and ensuring it matches the \texttt{prevHash} of the subsequent block. The second layer imports the Notary Authority's public key and verifies every signature. The engine reports the sequence number (\texttt{firsTamperedBlock}or \texttt{firstForgedSignatureBlock}) where any integrity check fails. The ledger components also contains the attacker simulator \texttt{attacks.js}, that provides a naive attack that only tamper the last digest and a more complex approach where, after a digest alteration, the entire hash chain is recalculated, which make the first  layer verification passes without error.

\begin{lstlisting}[language=JavaScript, caption=ledger/auditLog.js]
export async function verifyReceiptSignature(receipt) {...}

export async function auditLedger(entries) {
  if (entries.length === 0) return { status: "EMPTY", totalBlocks: 0 };

  let chainIntegrityPassed = true;
  let signatureIntegrityPassed = true;
  let firstTamperedBlock = -1;
  let firstForgedSignatureBlock = -1;
  let expectedPrevHash = GENESIS_HASH;

  for (let i = 0; i < entries.length; i++) {
    const block = entries[i];
    const recalculated = await computeChainHash(
      block.chainProof.prevHash,
      block.signature,
      block.docDigest,
    );

    const linkBroken =
      block.chainProof.prevHash !== expectedPrevHash ||
      recalculated !== block.chainProof.chainHash;
    if (linkBroken && chainIntegrityPassed) {
      chainIntegrityPassed = false;
      firstTamperedBlock = i;
    }
    expectedPrevHash = block.chainProof.chainHash;

    if (!(await verifyReceiptSignature(block)) && signatureIntegrityPassed) {
      signatureIntegrityPassed = false;
      firstForgedSignatureBlock = i;
    }
  }

  return {
    status: "OK",
    totalBlocks: entries.length,
    chainIntegrityPassed,
    signatureIntegrityPassed,
    firstTamperedBlock,
    firstForgedSignatureBlock,
  };
}

\end{lstlisting}


\subsection{Traffic generator}
The synthetic workload is managed by \texttt{poissonGenerator.js}. It uses exponential inter-arrival times derived via the inverse transform method:  $$\text{delay} = (-\frac{\ln(1-\text{rand}())} {\lambda}) \times 1000 \text{ milliseconds}.$$ A recursive timeout function schedules subsequent submissions, allowing to continously simulate requests matching the user-configured arrival late ($\lambda$)


\begin{lstlisting}[language=JavaScript, caption=sim/posissonGenerator.js]
export function poissonDelayMs(lambda, rand = Math.random) {...}

export function createPoissonGenerator({ submit, onStateChange = () => {}, log = () => {} }) {
  let active = false;

  function stop() {
    if (!active) return;
    active = false;
    onStateChange(false);
  }

  function start({ lambda, total, algorithm = "ECDSA" }) {
    active = true;
    onStateChange(true);
    let count = 0;
    log(`🚀 Poisson generator started (λ=${lambda} req/s, target=${total})`);

    function scheduleNext() {
      if (!active) return;
      if (count >= total) {
        stop();
        log(`🏁 Finished generating ${total} Poisson requests.`);
        return;
      }
      setTimeout(async () => {
        if (!active) return;
        count++;
        const digest = await sha256Bytes(`Poisson doc #${count} at ${performance.now()}`);
        submit(digest, algorithm);
        scheduleNext();
      }, poissonDelayMs(lambda));
    }

    scheduleNext();
  }

  return { start, stop, isActive: () => active };
}

\end{lstlisting}




\pagebreak


\section{Experimental Evaluation}

\subsection{Setup}
To evaluate cryptographic performance and queueing behavior without cross-run carryover, four independent workloads were executed in the browser (Brave) after refreshing the page prior to each test, ensuring an empty audit ledger and an unpopulated $100$-sample telemetry window at the start of every run:
\begin{enumerate}
    \item \textbf{ECDSA, low load:} $N = 100$ requests at configured $\lambda_{\text{cfg}} = 10\text{ req/s}$.
    \item \textbf{ECDSA, high load:} $N = 1000$ requests at configured $\lambda_{\text{cfg}} = 500\text{ req/s}$.
    \item \textbf{RSA-PSS, low load:} $N = 100$ requests at configured $\lambda_{\text{cfg}} = 10\text{ req/s}$.
    \item \textbf{RSA-PSS, high load:} $N = 1000$ requests at configured $\lambda_{\text{cfg}} = 500\text{ req/s}$.
\end{enumerate}

\vspace{1em}\noindent Because \texttt{computeQueueStatistics} requires at least five completions before outputting metrics, CSV telemetry begins at sequence index $4$ in every run (spanning indices $4\text{--}99$ for $N=100$ and $4\text{--}999$ for $N=1000$). Service time $S$ was timed strictly around signing and hash chaining inside the Web Worker via \texttt{performance.now()} ($\approx 0.1\text{ ms}$ timer resolution), while waiting time $W_q = t_{\text{QueueExit}} - t_{\text{Arrival}}$ was recorded by the main-thread FIFO queue. Analytical predictions ($\rho$, M/M/1, and M/G/1) were evaluated over the sliding window of the latest $100$ samples using $\lambda_{\text{cfg}}$.









\subsection{Service time of ECDSA vs RSA-PSS}
Table~\ref{tab:eval-summary} summarizes the final $100$-sample sliding-window telemetry recorded at the completion of each run.

\begin{table}[htbp]
\centering
\footnotesize
\setlength{\tabcolsep}{3.5pt}
\caption{Summary of final $100$-sample sliding-window telemetry for ECDSA (P-256) and RSA-PSS (4096-bit), comparing configured ($\lambda_{\text{cfg}}$) and measured ($\lambda_{\text{act}}$) arrival rates.}
\label{tab:eval-summary}
\begin{tabular}{lcccccccc}
\toprule
Algorithm & $\lambda_{\text{cfg}}$ & $\lambda_{\text{act}}$ & $\bar{S}$ & $\mu$ & $\rho$ & Emp.\ $\bar{W}_q$ & \multicolumn{2}{c}{Model $\bar{W}_q$ (ms)} \\
\cmidrule(lr){8-9}
 & (req/s) & (req/s) & (ms) & (req/s) & ($\lambda_{\text{cfg}}/\mu$) & (ms) & M/M/1 & M/G/1 \\
\midrule
ECDSA   & $10$  & $9.57$  & $0.160$ & $6{,}250$ & $0.0016$ & $0.052$   & $<0.001$ & $<0.001$ \\
ECDSA   & $500$ & $22.36$ & $0.133$ & $7{,}519$ & $0.0665$ & $0.018$   & $0.009$  & $0.006$ \\
RSA-PSS & $10$  & $9.90$  & $3.610$ & $277.0$   & $0.0361$ & $0.061$   & $0.135$  & $0.069$ \\
RSA-PSS & $500$ & $21.66$ & $3.576$ & $279.6$   & $1.7880$ & $275.236$ & $\infty$ & $\infty$ \\
\bottomrule
\end{tabular}
\end{table}

\noindent
The measurements demonstrate a clear separation in cryptographic cost between the two algorithms. ECDSA (P-256) achieves a mean service time of $\bar{S} = 0.160\text{ ms}$ ($\sigma_S \approx 0.077\text{ ms}$ over Seq.\ $4\text{--}99$, $\mu = 6{,}250\text{ req/s}$) at $\lambda_{\text{cfg}} = 10\text{ req/s}$, and $\bar{S} = 0.133\text{ ms}$ ($\mu = 7{,}519\text{ req/s}$) at $\lambda_{\text{cfg}} = 500\text{ req/s}$. Individual ECDSA readings lie near the $0.1\text{ ms}$ browser timer quantization floor, frequently registering $0.0\text{--}0.2\text{ ms}$.

\noindent
In contrast, 4096-bit RSA-PSS exhibits a mean service time of $\bar{S} = 3.610\text{ ms}$ ($\sigma_S \approx 0.532\text{ ms}$, $\mu = 277.0\text{ req/s}$) at $\lambda_{\text{cfg}} = 10\text{ req/s}$ and $\bar{S} = 3.576\text{ ms}$ ($\mu = 279.6\text{ req/s}$) at $\lambda_{\text{cfg}} = 500\text{ req/s}$. RSA-PSS signing is therefore approximately $22.6\times$ to $26.9\times$ slower than ECDSA, reflecting the heavy computational cost of 4096-bit modular exponentiation relative to P-256 elliptic-curve signing. Moreover, the RSA-PSS low-load run yields a very small squared coefficient of variation ($C_S^2 = \sigma_S^2 / \bar{S}^2 \approx 0.022 \ll 1$), confirming that cryptographic service times are nearly deterministic. Consequently, the Pollaczek--Khinchine M/G/1 formula predicts roughly half the waiting time of the exponential M/M/1 model ($0.069\text{ ms}$ vs.\ $0.135\text{ ms}$).

\subsection{Waiting time versus load}
Comparing waiting times across both arrival rates highlights the transition from lightly loaded operation to queue saturation:

\begin{itemize}
    \item \textbf{Light load ($\lambda_{\text{cfg}} = 10\text{ req/s}$):} With a fresh ledger, actual arrival rates closely matched the target rate ($9.57\text{ req/s}$ for ECDSA and $9.90\text{ req/s}$ for RSA-PSS). Utilization remained minimal ($\rho = 0.0016$ and $\rho = 0.0361$), so jobs rarely encountered a busy worker and empirical mean wait times stayed negligible at $0.052\text{ ms}$ and $0.061\text{ ms}$, respectively. Notably, for RSA-PSS, the empirical mean wait ($0.061\text{ ms}$) aligns closely with the M/G/1 prediction ($0.069\text{ ms}$), whereas M/M/1 ($0.135\text{ ms}$) overestimates delay by more than a factor of two.

    \item \textbf{ECDSA at high target load ($\lambda_{\text{cfg}} = 500\text{ req/s}$):} Because ECDSA's service capacity ($\mu = 7{,}519\text{ req/s}$) far exceeds the target arrival rate, configured utilization stayed well below saturation ($\rho = 0.0665$ at Seq.\ $999$). Across all $996$ logged requests, individual wait times never exceeded $0.1\text{ ms}$, and the final $100$-sample moving mean converged to $\bar{W}_q = 0.018\text{ ms}$, matching the sub-millisecond M/M/1 ($0.009\text{ ms}$) and M/G/1 ($0.006\text{ ms}$) estimates.

    \item \textbf{RSA-PSS at high target load ($\lambda_{\text{cfg}} = 500\text{ req/s}$):} With $\lambda_{\text{cfg}} = 500 > \mu \approx 279.6\text{ req/s}$, configured utilization remained strictly supercritical throughout the run ($\rho \in [1.640, 1.829]$, ending at $\rho = 1.7880$), causing both M/M/1 and M/G/1 models to diverge to $\infty$. Unlike ECDSA, the RSA-PSS queue never drained across the entire $1000$-request workload ($W_q > 0$ on every logged row): individual wait times grew from $9.8\text{ ms}$ at Seq.\ $4$ to a peak of $410.3\text{ ms}$ at Seq.\ $965$, driving the empirical $100$-sample mean wait $\bar{W}_q$ upward from $3.90\text{ ms}$ (Seq.\ $4$) to $275.236\text{ ms}$ at Seq.\ $999$ (peaking at $275.477\text{ ms}$ at Seq.\ $998$).
\end{itemize}

\subsection{Discussion}
The refreshed datasets confirm two key theoretical expectations while revealing important browser-level implementation effects at high request counts:

\begin{enumerate}
    \item \textbf{Validity of the M/G/1 model:} Under clean sub-critical traffic ($\lambda_{\text{cfg}} = 10\text{ req/s}$, RSA-PSS), empirical waiting times validate the M/G/1 Pollaczek--Khinchine prediction ($0.061\text{ ms}$ empirical vs.\ $0.069\text{ ms}$ M/G/1), confirming that low-variance cryptographic execution deviates substantially from the memoryless service assumption of M/M/1 ($0.135\text{ ms}$).

    \item \textbf{Browser timer clamping and effective arrival rate:} At $\lambda_{\text{cfg}} = 500\text{ req/s}$, initial inter-arrival intervals at Seq.\ $4\text{--}10$ were $4.1\text{--}6.5\text{ ms}$ rather than the theoretical $2.0\text{ ms}$ mean because browsers clamp nested \texttt{setTimeout} callbacks in \texttt{poissonGenerator.js} to a $4\text{ ms}$ minimum ($\le 250\text{ req/s}$). Over the full $1000$-request runs, average measured throughput dropped further to $22.36\text{ req/s}$ (ECDSA) and $21.66\text{ req/s}$ (RSA-PSS).

    \item \textbf{Main-thread DOM rendering overhead at high $N$:} Inspection of arrival timestamps shows inter-event intervals expanding steadily from $\approx 4\text{--}6\text{ ms}$ at Seq.\ $5\text{--}15$ to $\approx 80\text{--}105\text{ ms}$ near Seq.\ $999$. This slowdown stems from \texttt{auditTable.render(publicAuditLog)} clearing and reconstructing all accumulated HTML rows on every \texttt{SIGN\_SUCCESS} before invoking \texttt{server.complete()}. In the ECDSA run, \texttt{setTimeout} was delayed alongside rendering, keeping the queue empty ($W_q \le 0.1\text{ ms}$). In the RSA-PSS run, however, jobs queued during the initial rapid burst never cleared as main-thread render latency accumulated on top of the $3.58\text{ ms}$ signing time, creating a persistent multi-request backlog that drove waiting times past $400\text{ ms}$.
\end{enumerate}

Overall, resetting state between runs successfully eliminated carryover bias and produced accurate low-load benchmarks. To isolate pure cryptographic queueing behavior at high request counts ($N \ge 300$), the interface should append new audit rows incrementally (or cap displayed DOM rows) rather than re-rendering the entire table on every completion.
\pagebreak
\subsection{Discussion}
Inspecting the arrival timestamps in the CSV logs uncovers two critical implementation and browser effects that explain why the system did not exhibit unbounded queue growth at $\rho \approx 1.78$ and why empirical wait times exceeded analytical worker predictions during high-volume runs:

\begin{enumerate}
    \item \textbf{Main-thread DOM rendering bottleneck ($O(N^2)$ cumulative overhead):} In \texttt{main.js}, every \texttt{SIGN\_SUCCESS} handler calls \texttt{auditTable.render(publicAuditLog)} \emph{before} calling \texttt{server.complete()}. Because \texttt{auditTable.js} clears and rebuilds the entire HTML \texttt{<tbody>} for all accumulated receipts on every completion, rendering $1{,}100\text{--}2{,}200$ table rows blocks the main thread for roughly $80\text{--}200\text{ ms}$ per request. This holds the FIFO server in the \texttt{PROCESSING} state long after the worker has finished its $3.5\text{ ms}$ cryptographic operation, inflating observed $W_q$ to $80\text{--}265\text{ ms}$ while simultaneously starving \texttt{poissonGenerator.js} timeouts.

    \item \textbf{Configured vs.\ actual arrival rate and timer clamping:} Because recursive \texttt{setTimeout} calls in browsers are clamped to a minimum $4\text{ ms}$ delay ($\le 250\text{ req/s}$) and further delayed by main-thread rendering stalls, the \emph{actual} measured arrival rates were dramatically lower than the configured slider values: $18.88\text{ req/s}$ (ECDSA, $\lambda_{\text{cfg}} = 500$), $7.56\text{ req/s}$ (RSA-PSS, $\lambda_{\text{cfg}} = 500$), $4.75\text{ req/s}$ (RSA-PSS, $\lambda_{\text{cfg}} = 10$), and $4.54\text{ req/s}$ (ECDSA, $\lambda_{\text{cfg}} = 10$). Because \texttt{computeQueueStatistics} calculates $\rho = \lambda_{\text{cfg}} / \mu$ rather than using measured inter-arrival times, the theoretical models assumed $\lambda = 500\text{ req/s}$ ($\rho > 1.7$), whereas the worker's true utilization was $\lambda_{\text{actual}} / \mu \approx 0.027$. This explains why the queue periodically drained and why $\bar{W}_q$ stayed bounded around $58\text{--}161\text{ ms}$.

    \item \textbf{Sliding-window carryover across sequential runs:} Because the $100$-sample telemetry window is not cleared when starting a new workload, sequential tests without a page refresh retain carryover values in the first $99$ rows (visible at Seq.\ $1100$, $2100$, and $2200$).
\end{enumerate}

In summary, the experiments validate that ECDSA P-256 outperforms 4096-bit RSA-PSS by over $20\times$ in pure worker signing time and confirm that cryptographic service times have low variance ($C_S^2 \approx 0.04$), favoring M/G/1 over M/M/1. At scale, however, overall system latency in this architecture is governed not by cryptographic computation, but by $O(N)$ per-block DOM table re-rendering on the main thread---demonstrating that production audit interfaces should append rows incrementally rather than re-rendering the entire ledger.

\pagebreak

% ============================== 7. CONCLUSION (~0.5 page) ===================
\section{Conclusion}
\todo{Summarise results in 3--4 sentences, state one limitation, and give two
concrete future-work items (e.g.\ multiple workers $\to$ M/G/c, persistent log,
external timestamping).}

\end{spacing}

% ============================== REFERENCES ==================================
\begin{thebibliography}{9}
\bibitem{haber}  S.~Haber, W.~S.~Stornetta. How to time-stamp a digital document.
  \emph{Journal of Cryptology}, 3(2), 1991.
\bibitem{fips}   NIST. \emph{Digital Signature Standard (DSS)}, FIPS 186-5, 2023.
\bibitem{rfc}    K.~Moriarty et al. \emph{PKCS \#1: RSA Cryptography Specifications
  Version 2.2}, RFC 8017, 2016.
\bibitem{kleinrock} L.~Kleinrock. \emph{Queueing Systems, Volume 1: Theory}. Wiley, 1975.
\bibitem{webcrypto} W3C. \emph{Web Cryptography API}. \url{https://www.w3.org/TR/WebCryptoAPI/}
\end{thebibliography}

\end{document}
